%IEEE Access submission manuscript (template: ACCESS_latex_template_20260513)
%Ported from the Springer Nature sn-jnl version, content unchanged.
\documentclass{ieeeaccess}
\usepackage{cite}
\usepackage{amsmath,amssymb,amsfonts}
\usepackage{graphicx}
\usepackage{textcomp}
\usepackage{bm}
\usepackage{multirow}
\usepackage{booktabs}
\usepackage{algorithm}
\usepackage{algorithmicx}
\usepackage{algpseudocode}
\usepackage{url}

% Float placement tuning: let two wide figures share a page top and pack float pages
% tightly, so figures land near their first reference instead of trailing past it.
\setcounter{dbltopnumber}{3}
\setcounter{totalnumber}{5}
\renewcommand{\topfraction}{0.95}
\renewcommand{\dbltopfraction}{0.95}
\renewcommand{\textfraction}{0.05}
\renewcommand{\floatpagefraction}{0.8}
\renewcommand{\dblfloatpagefraction}{0.8}

% Clickable, colored cross-references and citations
\usepackage[colorlinks=true,linkcolor=blue,citecolor=blue,urlcolor=blue]{hyperref}

\makeatletter
\AtBeginDocument{\DeclareMathVersion{bold}
\SetSymbolFont{operators}{bold}{T1}{times}{b}{n}
\SetSymbolFont{NewLetters}{bold}{T1}{times}{b}{it}
\SetMathAlphabet{\mathrm}{bold}{T1}{times}{b}{n}
\SetMathAlphabet{\mathit}{bold}{T1}{times}{b}{it}
\SetMathAlphabet{\mathbf}{bold}{T1}{times}{b}{n}
\SetMathAlphabet{\mathtt}{bold}{OT1}{pcr}{b}{n}
\SetSymbolFont{symbols}{bold}{OMS}{cmsy}{b}{n}
\renewcommand\boldmath{\@nomath\boldmath\mathversion{bold}}}
\makeatother

\def\BibTeX{{\rm B\kern-.05em{\sc i\kern-.025em b}\kern-.08em
    T\kern-.1667em\lower.7ex\hbox{E}\kern-.125emX}}

\begin{document}
\history{Date of publication xxxx 00, 0000, date of current version xxxx 00, 0000.}
\doi{10.1109/ACCESS.2024.0429000}

\title{Entropy-Gated Contrastive Decoding: Mitigating Visual Inertia and Multimodal Hallucinations in Vision-Language Models}
\author{\uppercase{First A. Author}\authorrefmark{1}, \uppercase{Second B. Author}\authorrefmark{1}, AND \uppercase{Third C. Author}\authorrefmark{1}}

\address[1]{Department of Computer Science, Affiliation Name, Street Address, City 000000, State, Country (e-mail: corresponding.author@example.org; second.author@example.org; third.author@example.org).}

\markboth
{Author \headeretal: Entropy-Gated Contrastive Decoding for Vision-Language Models}
{Author \headeretal: Entropy-Gated Contrastive Decoding for Vision-Language Models}

\corresp{Corresponding author: First A. Author (e-mail: corresponding.author@example.org).}

\begin{abstract}
Vision-Language Models answer relational questions about objects well below their recognition accuracy. On the relation subset of the AMBER benchmark, most mainstream systems score near 0.7 accuracy, and the errors follow a measurable pattern: cross-attention settles on the earliest salient objects and pre-trained language priors finish the answer, a pattern known as Visual Inertia. We hypothesise that it underlies these errors and target the decoding-time decision point at which it would first appear. We present Entropy-Gated Contrastive Decoding (EGCD), a routing framework with frozen weights and per-model threshold calibration on an 822-query tuning split, that decides per query. Low entropy means the model has already resolved the query, and standard greedy decoding proceeds. A high-entropy first-token distribution routes the query to a counterfactual contrastive pathway: a second forward pass omits the image, an Adaptive Plausibility Constraint bounds the contrast to plausible tokens, and the isolated language prior is subtracted from the expert logits. EGCD raises the biased architectures over their greedy baselines: Qwen2-VL-2B from 70.71\% to 80.45\% balanced accuracy, Phi-3.5-Vision-4.2B from 69.05\% to 77.11\%, and SmolVLM2-2.2B from 63.51\% to 67.33\%. On the naturally balanced LLaVA models, where universal contrastive decoding collapses balanced accuracy by up to 10.62 points, the gate stays idle and returns the baseline essentially unchanged. Gating on predictive uncertainty turns contrastive decoding into a conditional, per-model calibrated intervention: applied where it helps, withheld where it harms.
\end{abstract}

\begin{keywords}
contrastive decoding, entropy-based routing, multimodal hallucination, vision-language models, visual inertia
\end{keywords}

\titlepgskip=-21pt

\maketitle

\section{Introduction}\label{sec:intro}

\PARstart{L}{arge} Vision-Language Models integrate visual perception with autoregressive text generation \cite{liu2023visual} and caption images, answer questions, and follow multimodal instructions with high fluency. Their record on fine-grained spatial and relational reasoning is far weaker. On the relation subset of AMBER, most mainstream models score near 0.7 accuracy on questions that ask whether two objects are in direct contact, and even the strongest systems misjudge whether objects are in contact \cite{wang2023amber}. Perturbation makes the same point from a different angle. Rotating images or adding noise brings down relational accuracy across open and closed models, but object recognition barely changes \cite{shin2026when}. Layout probes find the same pattern: when the structure is compositional, models fall back on non-visual shortcuts instead of actually reading the scene \cite{brown2025hallucination}. These results point to failures that arise at the decoding stage rather than at the input. Cross-attention locks onto the earliest salient objects and stays there, a phenomenon that \cite{gong2026attention} call Visual Inertia (Section~\ref{sec:related}), and that places the failure inside decoding rather than at the input.

The standard tool against distribution-level hallucination is Contrastive Decoding, which contrasts an expert distribution with a weaker amateur and subtracts the difference \cite{leng2024mitigating}. Applied uniformly, it has a documented failure mode (Section~\ref{sec:related}): much of its benchmark gain can come from a yes-ward shift of the output distribution rather than better grounding. Hidden-state interventions steer activations directly and avoid the logit arithmetic, but they can entangle unrelated features and disturb the natural generation behaviour of the model \cite{su2025activation,zhang2026mitigating}. What is missing is a way to apply the correction only where the failure actually occurs, and to leave every other query untouched.

We propose Entropy-Gated Contrastive Decoding (EGCD), a routing framework with frozen weights and per-model threshold calibration, built around one observation: the first-token distribution already separates queries the model answers confidently from those it does not. The router reads the Shannon entropy of that distribution. Low entropy means the model is confident and its priors are safe to trust. A sharp entropy spike activates a counterfactual contrastive pass that isolates the language prior and subtracts it, bounded by an Adaptive Plausibility Constraint. The same gate that helps bias-affected architectures stays idle on naturally balanced ones, which no uniform scheme can do. Our contributions are as follows:

\begin{enumerate}
\item A decoding-time uncertainty signal motivated by Visual Inertia \cite{gong2026attention}. We take their attention-stagnation account and reduce it to a measurable probe based on the first-token distribution, where entropy separates correct from incorrect answers on every model (binary error-prediction AUC between 0.62 and 0.79; model-level values in Section~\ref{subsec:dynamics}). This motivates intervention at decoding time.

\item EGCD, a routing mechanism that uses first-token Shannon entropy as a real-time uncertainty signal and triggers counterfactual contrastive decoding only when entropy crosses a calibrated threshold. The gate adds no parameters, requires no sampling, and reads a statistic that the standard pass already computes. The threshold $\tau$ is calibrated on a disjoint 822-query tuning split, with no weight updates.

\item A cross-architecture study on AMBER showing that contrastive interventions should be conditional: EGCD lifts balanced accuracy by up to 9.74 points on biased architectures and stays idle on naturally balanced ones, avoiding collapses of up to 10.62 points that universal contrastive decoding inflicts on them.
\end{enumerate}

The remainder of this paper is organised as follows. Section~\ref{sec:related} reviews related work on visual inertia, contrastive decoding, and uncertainty-based routing. Section~\ref{sec:method} introduces our proposed methodology, detailing the entropy gate and the counterfactual contrastive pathway. Section~\ref{sec:experiments} describes the experimental setup, including the AMBER benchmark, models, and baselines. Section~\ref{sec:results} presents the results and analysis across seven vision-language models. Finally, Section~\ref{sec:conclusion} concludes the paper and outlines limitations and future work.

\section{Related Work}\label{sec:related}

\subsection{Visual Inertia and Relational Hallucinations}\label{subsec:visual-inertia}

Hallucination benchmarks have moved well beyond simple object-existence checks. POPE samples object presence in a balanced way \cite{li2023bevaluating}; CHAIR counts invented objects in captions \cite{rohrbach2018object}; MME spans perceptual and cognitive tasks \cite{fu2025mme}; and AMBER adds a discriminative relation subset of 1,664 image-text pairs \cite{wang2023amber}. Relational queries are the hardest of the lot. A model has to bind two objects and then check a predicate between them, and that binding step is exactly what attention stagnation disrupts. The evidence comes from more than one direction. Rotation and noise take down relational accuracy while recognition survives \cite{shin2026when}; spatial benchmarks keep turning up shortcut heuristics that skip the grounded binding step \cite{brown2025hallucination}; and surveys of the field still list compositional failure as an open problem for vision-language modelling \cite{bordes2024introvlm}.

\cite{gong2026attention} call this mechanism Visual Inertia: attention to semantically critical regions settles early and then stays fixed, so the model never inspects the secondary objects that a relational query targets. Their counterfactual experiment is straightforward. An inertia-inducing intervention reduces visual activeness from 41.52 to 19.24 and the relation score from 63.96 to 47.25, which is the largest decline among their controls. Two consequences follow for intervention design. First, uniform attention amplification cannot help, because scaling a settled map only reinforces the settlement \cite{gong2026attention,chen2026mitigating}. Second, once the visual signal freezes, unimodal language priors dominate the deeper layers; this dominance has now been traced to specific FFN modules \cite{guo2026fade} and is the target of prior-suppression decoding methods \cite{ren2026nolan,min2025mitigating}. EGCD is designed around both of these points: it avoids uniform attention amplification and directly counteracts the language-prior dominance at decoding time.

\subsection{Contrastive Decoding and Hidden-State Interventions}\label{subsec:contrastive}

Contrastive decoding was first introduced for open-ended text generation, where an expert distribution is contrasted with the output of a weaker amateur model \cite{li2023acontrastive,su2022empirical}. Visual Contrastive Decoding (VCD) adapted this idea to hallucination research by using a Gaussian-noised image as the amateur, with an Adaptive Plausibility Constraint (APC) that restricts the contrast to plausible tokens \cite{leng2024mitigating}. Subsequent variants have replaced the amateur with instruction distractors \cite{wang2024mitigating}, context-and-text-aware token selection \cite{huo2025self}, counterfactual masks \cite{xiao2026macd}, or an omitted visual prompt \cite{favero2024multi}. More recent work modulates the contrast adaptively based on query type or entropy \cite{im2025selfaug,huang2026chasd}.

The uniform application of this recipe has a documented failure mode. \cite{yin2025mirage} show that the reported POPE gains of contrastive methods come largely from a yes-ward shift in the output distribution, and that under sampling, the APC degrades decoding into an approximation of greedy search; a spurious prompt adjustment can reproduce the same numbers. A reproducibility study by \cite{bendre2026rethinking} extends this critique across multiple models and decoding regimes. Hidden-state interventions take a different approach by editing activations rather than logits, including bidirectional activation steering \cite{su2025activation} and selective interventions that attempt to leave the token distribution intact \cite{zhang2026mitigating}, an idea also explored for contextual faithfulness in text-only language models \cite{wang2025contextfocus}; these operations are distinct from layer-logit contrast methods such as DoLa \cite{chuang2024dola}, which read and contrast logits at intermediate layers rather than editing the activations themselves. Steering signals can entangle unrelated features and shift the model away from its natural generation behaviour. EGCD operates at the logit level, runs only when the gate fires, and decodes greedily, which closes off the sampling-collapse route to spurious gains by construction. Section~\ref{sec:results} evaluates both under one protocol and shows that where each method helps, harms, or does nothing depends on the amateur and on the model.

\subsection{Uncertainty Estimation and Dual-Process Routing}\label{subsec:uncertainty}

Dual-process theories split cognition into a fast, automatic mode and a slower, deliberate one \cite{wason1974dual,kahneman2011thinking}, and the same split now runs through reasoning research on LLMs \cite{li2025system}. DynaThink, for instance, routes a question to fast thinking only when sampled chains agree and the majority answer needs the fewest reasoning steps \cite{pan2024dynathink}; HaluSearch interleaves fast and slow search to cut hallucination \cite{cheng2025think}. The catch is that both routers verify by sampling. Computing the routing signal costs more than producing the answer it guards.

Octopus is the closest prior work to ours. It learns a per-step router that picks among contrastive variants through a trained decision token \cite{suo2025octopus}. The learned router covers more cases, but it needs a preference-optimisation stage and pays overhead at every token. Probing studies suggest cheaper routes: linear probes pull structured information out of intermediate representations \cite{alain2016understanding}, and tuned lenses recover latent predictions layer by layer \cite{belrose2023eliciting}. EGCD uses none of this: it routes once per query, on a first-token statistic the standard pass already produces.

\Figure[t!p](topskip=0pt, botskip=0pt, midskip=0pt)[width=\textwidth]{fig1.png}{Overview of Entropy-Gated Contrastive Decoding (schematic). A relational query enters the frozen VLM; the first-token entropy $H(Y_{1})$ is compared against $\tau$, and the query is routed to System 1 (greedy decoding) or System 2 (counterfactual contrastive decoding with APC; $\alpha = 1.0$, $\beta = 0.1$). The selected-threshold interval across the seven evaluated models is $\tau^{*} \in [0.570, 0.700]$; the search grid is $[0.550, 0.750]$.\label{fig:overview}}

\section{Methodology}\label{sec:method}

This section formalizes EGCD: the notation, the gate, and the corrective pathway. Figure~\ref{fig:overview} traces a relational query through both routes to the final response.

\subsection{Problem Formulation and Visual Inertia}\label{subsec:problem}

Let $\theta$ denote the parameters of a model, frozen at inference time, and let $V$ denote its vocabulary. The input is a tuple $(X_{v}, X_{t})$ of a visual scene and a text prompt, and autoregressive decoding produces $Y = (y_{1}, \ldots, y_{T})$. At step $t$ the standard expert multimodal pass produces a vocabulary logit vector $z_{\text{std}} \in \mathbb{R}^{|V|}$ whose component for token $y$ we write as $z_{\text{std}}(y)$. The standard next-token distribution is
\begin{equation}\label{eq:std}
P_{\text{std}}(y \mid X_{v}, X_{t}, Y_{<t}) \;=\; \frac{\exp\, z_{\text{std}}(y)}{\sum_{w \in V} \exp\, z_{\text{std}}(w)}.
\end{equation}

Visual Inertia \cite{gong2026attention} is the failure mode we target. The spatial cross-attention map stays anchored to early salient objects as decoding proceeds, so the grounding signal decays and the statistical language prior takes over the next-token distribution. In binary relation tasks the decoder defaults to conservative answers and stereotypical co-occurrences \cite{yin2025mirage}. The substitution is measurable: it shows up as probability placed on answer tokens that the image does not support. A correction has to fire exactly when this substitution happens, and stay quiet otherwise.

\subsection{Query-Level Entropy Gating and Dual-Process Routing}\label{subsec:gate}

Penalising logits at every step taxes queries that need no help, so we treat decoding as a dual-process system \cite{kahneman2011thinking} and route once, at the query level. The probe is the Shannon entropy \cite{shannon1948mathematical} of the first-token distribution over the full vocabulary $V$:
\begin{equation}\label{eq:entropy}
H(Y_{1}) \;=\; -\sum_{y \in V} P_{\text{std}}(y \mid X_{v}, X_{t})\, \log P_{\text{std}}(y \mid X_{v}, X_{t}).
\end{equation}
$H(Y_{1})$ is the full-vocabulary Shannon entropy in nats, computed from the unscaled logits of the standard expert pass. Its maximum value is $\log |V|$, so the threshold grid $[0.550, 0.750]$ used in Section~\ref{subsec:baselines} is well within the supported range. Throughout, ``temperature zero'' refers to greedy token selection at generation time; the entropy itself is always computed from the unmodified softmax distribution.

A peaky distribution yields low entropy. When language priors and visual evidence diverge at the relational decision point, probability mass diffuses across candidates and entropy spikes. For the binary relation queries used in this paper, the prompt is constructed so that the first generated token is the answer itself; the chat template, leading-whitespace handling, and Yes/No tokenisation are described in Section~\ref{subsec:baselines}.

A single threshold $\tau$ separates the two cases. Queries with $H(Y_{1}) < \tau$ are treated as settled: their priors can be trusted, and System 1 answers them with standard greedy decoding, giving $Y_{\text{std}}$. When $H(Y_{1}) \geq \tau$, the query goes to the corrective System 2 pathway, which returns $Y_{\text{CD}}$. The gate reads a forward-pass statistic that System 1 already produces, so the routing check itself adds no forward pass; the extra cost (one text-only pass for the first token of a routed query; for longer outputs, one standard and one counterfactual pass per additional token) falls only on queries that actually route to System 2.

\subsection{Counterfactual Isolation and Adaptive Plausibility Constraint}\label{subsec:apc}

System 2 isolates the language prior through a counterfactual forward pass that removes the visual input entirely \cite{favero2024multi}:
\begin{equation}\label{eq:cf}
z_{\text{cf}}(y) \;=\; \text{logit}_{\theta}\bigl(y \mid X_{t}, Y_{<t}\bigr),
\end{equation}
i.e.\ $z_{\text{cf}}$ is the text-only conditional logit vector at decoding step $t$, used as a proxy for the language prior. The contrastive subtraction follows the same operationalisation as VCD \cite{leng2024mitigating} and counterfactual visual grounding \cite{favero2024multi}.

Subtracting $z_{\text{cf}}$ raw can amplify tail tokens, because a strongly negative prior logit becomes a large positive score after subtraction, which invites ungrammatical output. Following VCD \cite{leng2024mitigating}, we bound the contrast to an adaptive plausibility set. With the context $Y_{<t}$ and inputs $(X_{v}, X_{t})$ suppressed for readability, the set is
\begin{align}\label{eq:head}
V_{\text{head}} \;=\;& \bigl\{\, y \in V \,:\, P_{\text{std}}(y) \;\geq\; \beta\, \max_{w \in V} P_{\text{std}}(w)\,\bigr\},\notag\\
& 0 < \beta \leq 1.
\end{align}

For every plausible token, the contrastive logits combine the expert and counterfactual terms; all other tokens are truncated to negative infinity:
\begin{align}\label{eq:cd}
z_{\text{CD}}(y) \;=\;& (1 + \alpha)\, z_{\text{std}}(y) \;-\; \alpha\, z_{\text{cf}}(y),\notag\\
& \forall\, y \in V_{\text{head}}, \quad \alpha \geq 0.
\end{align}

The calibrated distribution is the softmax over this truncated space (equivalently, proportional to $P_{\text{std}}(y)^{1+\alpha} / P_{\text{cf}}(y)^{\alpha}$ on $V_{\text{head}}$), and $Y_{\text{CD}}$ is decoded greedily. Because System 2 never samples, the truncation cannot manufacture greedy-search gains of the kind documented under sampling decoders \cite{yin2025mirage,bendre2026rethinking}; the contrast acts only where the gate routed the query. Algorithm~\ref{alg:egcd} summarises the full procedure.

\begin{algorithm}
\caption{Entropy-Gated Contrastive Decoding (EGCD)}\label{alg:egcd}
\begin{algorithmic}[1]
\Require Frozen VLM $p_{\theta}$ with logit function $z_{\theta}$
\Require Multimodal input $(X_{v}, X_{t})$
\Require Entropy threshold $\tau$, contrastive strength $\alpha \geq 0$, plausibility threshold $0 < \beta \leq 1$
\Require Maximum decoding length $T$
\State $Y \leftarrow \emptyset$; \; $y \leftarrow \langle\text{BOS}\rangle$ \Comment{shared prefix; current token}
\State $z_{\text{std}} \leftarrow z_{\theta}(\,\cdot \mid X_{v}, X_{t})$; \; $P_{\text{std}} \leftarrow \text{softmax}(z_{\text{std}})$ \Comment{gate pass; reused for the first decision in both branches}
\State $H(Y_{1}) \leftarrow -\sum_{y \in V} P_{\text{std}}(y)\, \log P_{\text{std}}(y)$
\If{$H(Y_{1}) < \tau$} \Comment{System 1: standard greedy decoding}
  \While{$y \neq \langle\text{EOS}\rangle$ \textbf{and} $|Y| < T$}
    \State $y \leftarrow \operatorname*{arg\,max}_{y'} P_{\text{std}}(y')$; \; append $y$ to $Y$
    \If{$y \neq \langle\text{EOS}\rangle$ \textbf{and} $|Y| < T$} \Comment{recompute only if another token follows}
      \State $z_{\text{std}} \leftarrow z_{\theta}(\,\cdot \mid X_{v}, X_{t}, Y)$; \; $P_{\text{std}} \leftarrow \text{softmax}(z_{\text{std}})$
    \EndIf
  \EndWhile
\Else \Comment{System 2: counterfactual contrastive decoding}
  \While{$y \neq \langle\text{EOS}\rangle$ \textbf{and} $|Y| < T$}
    \State $z_{\text{cf}} \leftarrow z_{\theta}(\,\cdot \mid X_{t}, Y)$ \Comment{text-only pass; $z_{\text{std}}$ reused from the gate pass on the first iteration}
    \State $V_{\text{head}} \leftarrow \{\, y \in V \,:\, P_{\text{std}}(y) \geq \beta \cdot \max_{w} P_{\text{std}}(w)\,\}$
    \State $z_{\text{CD}}(y) \leftarrow (1 + \alpha)\, z_{\text{std}}(y) - \alpha\, z_{\text{cf}}(y)$ \textbf{for} $y \in V_{\text{head}}$; $-\infty$ \textbf{otherwise}
    \State $y \leftarrow \operatorname*{arg\,max}_{y'} \text{softmax}(z_{\text{CD}})(y')$; \; append $y$ to $Y$
    \If{$y \neq \langle\text{EOS}\rangle$ \textbf{and} $|Y| < T$}
      \State $z_{\text{std}} \leftarrow z_{\theta}(\,\cdot \mid X_{v}, X_{t}, Y)$; \; $P_{\text{std}} \leftarrow \text{softmax}(z_{\text{std}})$
    \EndIf
  \EndWhile
\EndIf
\State \Return $Y$
\end{algorithmic}
\end{algorithm}

The shared prefix $Y$ is appended to in both branches, and the EOS check runs in both branches; System 2 maintains its own KV cache distinct from System 1, but the selected tokens are shared so that the two streams read from the same prefix. For the binary-relation evaluation in this paper, $T = 1$ (a single answer token), so the loop body executes once per routed query; longer generations would repeat the counterfactual pass at every token.

\section{Experimental Setup}\label{sec:experiments}

\subsection{Evaluation Benchmark}\label{subsec:benchmark}

We evaluate on the AMBER benchmark \cite{wang2023amber}, restricted to its Discriminative-Relation subset: 1,664 image-text pairs of binary Yes/No questions about direct physical contact between objects. Existence suites such as POPE \cite{li2023bevaluating} and captioning metrics such as CHAIR \cite{rohrbach2018object} rarely stress relational binding, and AMBER's authors report affirmative bias across mainstream models, with Yes-ratios exceeding No-ratios for nearly every evaluated system. The subset is balanced by construction at the dataset level. Relational verification requires shifting attention to secondary objects, which makes this subset a direct probe of the Visual Inertia hypothesis. Scoring well on this subset requires reading the relation out of the image, which is the property we need from a test bed.

\subsection{Target Vision-Language Models}\label{subsec:models}

The suite spans seven open-weight models from four architectural families. Qwen2-VL-2B and Qwen2.5-VL-7B use native dynamic resolution and multimodal rotary position embeddings \cite{bai2025qwen25vl}. SmolVLM2-2.2B continues the Idefics lineage in a compact form \cite{marafioti2025smolvlm,laurencon2024matters}, alongside Idefics2-8B itself \cite{laurencon2024matters}. LLaVA-1.5-7B and LLaVA-v1.6-Mistral-7B pair CLIP vision encoders with Vicuna and Mistral backbones through MLP projection layers \cite{liu2023visual,liu2024improved}. The last model in the suite is Phi-3.5-Vision-4.2B \cite{abdin2024phi3}. Encoders, connectors, and language backbones vary across all seven, which is deliberate: it lets us see how the architecture itself shapes the response to a decoding intervention.

\subsection{Evaluation Metrics}\label{subsec:metrics}

Recall is computed on the positive split of the benchmark, where the ground-truth relation holds and the expected answer is Yes. Specificity comes from the negative, hallucination-prone split, where the relation is false and an affirmative shortcut yields no gain \cite{yin2025mirage}. With True Positives defined as correct Yes answers to true relations, False Negatives as unjustified No answers to true relations, False Positives as hallucinated Yes answers to false relations, and True Negatives as correct No answers to false relations, the metrics are:
\begin{align}
\text{Recall} &= \frac{\text{True Positives}}{\text{True Positives} + \text{False Negatives}}, \label{eq:recall}\\
\text{Specificity} &= \frac{\text{True Negatives}}{\text{True Negatives} + \text{False Positives}}. \label{eq:spec}
\end{align}

Balanced Accuracy is the mean of Recall and Specificity \cite{brodersen2010balanced}, one term from each split:
\begin{equation}\label{eq:balacc}
\text{Balanced Accuracy} = \frac{\text{Recall} + \text{Specificity}}{2}.
\end{equation}

Together, recall and specificity characterise the bias profile of a method; balanced-accuracy gains reflect shifts in this profile rather than a direct causal test of visual grounding.

\subsection{Baseline Configurations and Implementation Constraints}\label{subsec:baselines}

All runs use zero-shot prompting with supervised threshold calibration on a disjoint tuning split; weights are frozen, half-precision, greedy decoding with temperature zero for reproducibility. EGCD uses $\alpha = 1.0$ and $\beta = 0.1$, the defaults established for visual contrastive decoding \cite{leng2024mitigating}. The gate threshold is the one design choice that varies per model.

We split the 1,664 relation pairs into a tuning half (822 pairs) and a held-out half (842 pairs) using a canonical partition stratified by image id, so that multi-query images never leak between partitions. The split is identical across all seven models and is released alongside the code. The selection objective on the tuning half is the argmax of balanced accuracy over the grid $\tau \in [0.550, 0.750]$ in steps of $0.01$; ties are broken deterministically by selecting the lowest $\tau$. Every result in the main text is reported on the held-out half, except the threshold sweep itself (Figures~\ref{fig:sweep}--\ref{fig:all} and Appendix~\ref{app:sweep}), which by construction reports tuning-half numbers. The held-out partition contains 483 positive and 359 negative queries; balanced accuracy rather than raw accuracy is therefore the headline metric.

The five decoding regimes compared per model are: the greedy baseline; universal CD+APC applied at every step; EGCD; VCD (Leng et al.\ 2024, Gaussian-noised image amateur with noise standard deviation $\sigma = 0.5$ and seed fixed for reproducibility); and a layer-logit contrast (LLC) baseline that hooks the decoder at layer 3 and projects through the LM head, following the implementation pattern of DoLa \cite{chuang2024dola}. We label this baseline ``LLC'' to distinguish it from Self-Introspective Decoding \cite{huo2025self}, which uses Context-and-Text-aware Token Selection (CT$^{2}$S), a different operation not implemented here. VCD and LLC both use $\alpha = 1.0$ and $\beta = 0.1$ when an adaptive plausibility constraint is applied; LLC reads the amateur logits from layer 3 and contrasts them with the final-layer expert logits.

The chat template, leading-whitespace handling, and Yes/No tokenisation are model-specific and follow each model's official inference code. For half-precision inference, log-softmax and entropy are computed in float32; a numerical floor of $10^{-12}$ is added before taking the logarithm to avoid NaNs from zero-probability tokens.

\section{Results and Analysis}\label{sec:results}

The evaluation covers the 842 held-out relational pairs across seven models in zero-shot, half-precision settings. We report three axes: gains on biased architectures, the boundary of contrastive interventions, and the dynamics of the gate itself. Throughout, EGCD denotes the gated method at its per-model operating point $\tau^{*}$, universal CD+APC denotes the same contrastive arithmetic applied at every step of every query, VCD and LLC denote the two published baselines of Section~\ref{subsec:baselines} applied the same way, and Baseline denotes plain greedy decoding.

\subsection{Quantitative Relational Grounding Results}\label{subsec:quant}

Table~\ref{tab:main} summarises the five decoding regimes. Three patterns stand out. Throughout, \emph{biased} and \emph{balanced} describe observed behaviour under contrastive decoding rather than measured architectural properties. Biased models improve sharply under EGCD. Balanced models are damaged by universal CD+APC, in one case catastrophically. On the held-out split EGCD never falls more than 0.25 points below its baseline, and it never trails universal CD+APC by more than 0.57 points on any model where CD+APC improves on greedy decoding, while avoiding the catastrophic collapses on the LLaVA models.

\begin{table*}[!t]
\centering
\caption{Held-out results (842 queries) for the five decoding regimes. EGCD uses the per-model threshold $\tau^{*}$ selected on the tuning half; VCD and LLC are applied at every decoding step. VCD uses $\alpha = 1.0$ and $\beta = 0.1$ as in CD+APC; LLC reads amateur logits from layer 3 and contrasts them with the final-layer expert logits.}\label{tab:main}
\begin{tabular}{@{}llccc@{}}
\toprule
\textbf{Model} & \textbf{Regime} & \textbf{Recall (\%)} & \textbf{Specificity (\%)} & \textbf{Bal.\ Acc. (\%)} \\
\midrule
\multirow{5}{*}{\textbf{Qwen2-VL-2B}}
 & Greedy baseline            & 47.00 & 94.43 & 70.71 \\
 & Universal CD+APC           & 74.95 & 79.67 & 77.31 \\
 & \textbf{EGCD (ours) $\tau^{*} = 0.70$} & \textbf{72.88} & \textbf{88.02} & \textbf{80.45} \\
 & VCD                        & 48.45 & 93.87 & 71.16 \\
 & LLC                        & 47.00 & 94.43 & 70.71 \\
\midrule
\multirow{5}{*}{\textbf{Phi-3.5-Vision-4.2B}}
 & Greedy baseline            & 46.17 & 91.92 & 69.05 \\
 & Universal CD+APC           & 75.57 & 77.99 & 76.78 \\
 & \textbf{EGCD (ours) $\tau^{*} = 0.57$} & \textbf{72.88} & \textbf{81.34} & \textbf{77.11} \\
 & VCD                        & 41.61 & 91.09 & 66.35 \\
 & LLC                        & 42.03 & 93.31 & 67.67 \\
\midrule
\multirow{5}{*}{\textbf{SmolVLM2-2.2B}}
 & Greedy baseline            & 31.47 & 95.54 & 63.51 \\
 & Universal CD+APC           & 49.69 & 84.40 & 67.05 \\
 & \textbf{EGCD (ours) $\tau^{*} = 0.68$} & \textbf{48.03} & \textbf{86.63} & \textbf{67.33} \\
 & VCD                        & 53.62 & 81.89 & 67.76 \\
 & LLC                        & 31.47 & 95.54 & 63.51 \\
\midrule
\multirow{5}{*}{\textbf{Qwen2.5-VL-7B}}
 & Greedy baseline            & 59.01 & 98.89 & 78.95 \\
 & Universal CD+APC           & 67.29 & 95.82 & 81.55 \\
 & \textbf{EGCD (ours) $\tau^{*} = 0.64$} & \textbf{64.18} & \textbf{97.77} & \textbf{80.98} \\
 & VCD                        & 59.01 & 97.77 & 78.39 \\
 & LLC                        & 59.01 & 98.89 & 78.95 \\
\midrule
\multirow{5}{*}{\textbf{Idefics2-8B}}
 & Greedy baseline            & 76.19 & 88.86 & 82.52 \\
 & Universal CD+APC           & 75.57 & 87.74 & 81.66 \\
 & \textbf{EGCD (ours) $\tau^{*} = 0.66$} & \textbf{76.81} & \textbf{87.74} & \textbf{82.28} \\
 & VCD                        & 77.64 & 84.96 & 81.30 \\
 & LLC                        & 76.60 & 88.86 & 82.73 \\
\midrule
\multirow{5}{*}{\textbf{LLaVA-1.5-7B}}
 & Greedy baseline            & 67.49 & 82.73 & 75.11 \\
 & Universal CD+APC           & 35.20 & 94.99 & 65.09 \\
 & \textbf{EGCD (ours) $\tau^{*} = 0.70$} & \textbf{67.49} & \textbf{82.73} & \textbf{75.11} \\
 & VCD                        & 75.98 & 70.75 & 73.37 \\
 & LLC                        & 72.05 & 76.32 & 74.19 \\
\midrule
\multirow{5}{*}{\textbf{LLaVA-v1.6-Mistral-7B}}
 & Greedy baseline            & 78.47 & 93.04 & 85.75 \\
 & Universal CD+APC           & 63.35 & 86.91 & 75.13 \\
 & \textbf{EGCD (ours) $\tau^{*} = 0.70$} & \textbf{78.26} & \textbf{93.04} & \textbf{85.65} \\
 & VCD                        & 75.57 & 91.36 & 83.47 \\
 & LLC                        & 78.47 & 93.04 & 85.75 \\
\bottomrule
\end{tabular}
\end{table*}

\Figure[t!p](topskip=0pt, botskip=0pt, midskip=0pt)[width=\textwidth]{fig2.png}{Balanced accuracy on the held-out half (842 queries) for the seven models under greedy decoding, Universal CD+APC, VCD, LLC, and EGCD, ordered by EGCD's gain over greedy decoding.\label{fig:balacc}}

Figure~\ref{fig:balacc} shows the balanced-accuracy comparison. On Qwen2-VL-2B, EGCD adds 9.74 points over greedy decoding and 3.14 over universal CD+APC. That 3.14 gap is the only EGCD-over-CD improvement on the four biased models that holds up under an exact paired McNemar test \cite{dietterich1998approximate} on per-item correctness ($p = 0.002$, 30 discordant pairs favouring EGCD versus 10 favouring CD). On the remaining three biased architectures, Phi-3.5-Vision-4.2B, SmolVLM2-2.2B, and Qwen2.5-VL-7B, the gap to CD is not significant under the same test ($p = 1.00, 1.00, 0.24$ respectively). McNemar's test compares per-item accuracy, and the reported p-values are unadjusted for multiple comparisons; the balanced-accuracy gaps in Table~\ref{tab:main} are descriptive. On the two LLaVA models, EGCD significantly outperforms universal CD because the gate stays idle where universal CD catastrophically fails (LLaVA-1.5: $p < 0.001$, 160 discordant pairs favouring EGCD versus 48 favouring CD; LLaVA-v1.6: $p < 0.001$, 123 discordant pairs favouring EGCD versus 29 favouring CD). The VCD and LLC rows make the same point from a different angle. VCD uses a noised image as its amateur and is the strongest single method on SmolVLM2-2.2B (67.76, against EGCD's 67.33). LLC reads amateur logits from the model's own shallow decoder layer; it stays within 2.1 points of the greedy baseline on every model and changes no answer on Qwen2-VL-2B, where the headroom is largest, or on SmolVLM2-2.2B, Qwen2.5-VL-7B, and LLaVA-v1.6-Mistral-7B. A universal amateur helps when its bias lines up with the model's failure mode, and does nothing when it does not. The gate is what connects the corrective pathway to the queries that actually need it.

\subsection{Resolving Visual Inertia in Biased Architectures}\label{subsec:resolve}

SmolVLM2-2.2B and Qwen2-VL-2B start with baseline recalls of 31.47\% and 47.00\%, respectively. In both cases, cross-attention stagnates and language priors dominate, so the models answer No to relations that actually hold. Universal CD+APC corrects this bias by shifting the output distribution in the opposite direction. On Qwen2-VL-2B, Recall climbs to 74.95\%, but Specificity falls from 94.43\% to 79.67\%. The shift is one-sided: recall rises by 28 points while specificity falls by 15, so the model trades a No-bias for a Yes-bias. Figure~\ref{fig:delta} traces the net effect of this trade-off on balanced accuracy.

\Figure[t!p](topskip=0pt, botskip=0pt, midskip=0pt)[width=\textwidth]{fig3.png}{Change in balanced accuracy relative to greedy decoding, per model and method: Universal CD+APC buys its gains with a 10--11 point collapse on the two LLaVA models, while EGCD keeps the gains without the collapse.\label{fig:delta}}

The gate makes the recall--specificity trade less costly rather than removing it. At the selected threshold, Qwen2-VL-2B reaches 80.45\% balanced accuracy, with Recall at 72.88\% and Specificity at 88.02\%: Recall ends up about two points short of universal CD (74.95\%) while Specificity falls far less (CD drops to 79.67\%). SmolVLM2 peaks at its threshold with 67.33\%, improving recall at a smaller Specificity cost than universal CD+APC. Figure~\ref{fig:recspec} compares Recall and Specificity for the four bias-affected models under all five decoding regimes. EGCD still sacrifices some specificity relative to greedy decoding (Qwen2-VL: 94.43 $\rightarrow$ 88.02; Phi-3.5: 91.92 $\rightarrow$ 81.34), but the cost is smaller than what universal CD+APC pays for the same recall.

\Figure[t!p](topskip=0pt, botskip=0pt, midskip=0pt)[width=\textwidth]{fig4.png}{Recall and Specificity on the held-out half for the four bias-affected models, where Universal CD+APC buys its recall by spending specificity and EGCD recovers most of the recall at a fraction of the specificity cost.\label{fig:recspec}}

Figure~\ref{fig:sweep} sweeps $\tau$ over $[0.550, 0.750]$ for the four bias-affected architectures. Each curve rises above both reference lines, hits a peak ($\tau^{*} = 0.70, 0.57, 0.68$, and $0.64$ for Qwen2-VL-2B, Phi-3.5-Vision-4.2B, SmolVLM2-2.2B, and Qwen2.5-VL-7B), and then falls back toward baseline once $\tau$ moves past the entropy mass of genuine conflicts.

\Figure[t!p](topskip=0pt, botskip=0pt, midskip=0pt)[width=\textwidth]{fig5.png}{Balanced accuracy across the entropy-threshold sweep on the tuning half for the four bias-affected models; the star marks each model's selected threshold $\tau^{*}$.\label{fig:sweep}}

\subsection{The Architectural Boundary of Contrastive Interventions}\label{subsec:boundary}

Contrastive decoding is not a universal tool. LLaVA-1.5-7B and LLaVA-v1.6-Mistral-7B start with recalls of 67.49\% and 78.47\%, which indicates that their connectors produce visual tokens that already compete on equal terms with text. Universal CD+APC damages exactly these models: on LLaVA-1.5-7B, balanced accuracy falls to 65.09\% with Recall dropping to 35.20\%, and on LLaVA-v1.6-Mistral-7B it drops by 10.62 points. We refer to this as contrastive overcompensation. For a model that is already balanced, subtracting a vision-free amateur over-penalises the syntactic and structural priors that the model still relies on, and recall suffers as a result. The same arithmetic that improves Qwen2-VL-2B ends up removing corrections that a balanced model was already making correctly.

The gate handles this distinction without any explicit instruction. Balanced models maintain compact first-token distributions, entropy rarely crosses $\tau^{*}$, and the gate does not activate. As a result, EGCD preserves 75.11\% on LLaVA-1.5-7B and 85.65\% on LLaVA-v1.6-Mistral-7B, matching greedy decoding within 0.10 percentage points (the gate routes at most two of the 842 held-out queries at the selected $\tau^{*}$). Figure~\ref{fig:sound} shows the sweep for the naturally balanced architectures: Idefics2-8B does not differ from greedy decoding (exact paired McNemar $p = 1.00$) or from universal CD+APC ($p = 0.07$), and LLaVA-1.5-7B recovers its baseline at $\tau^{*} = 0.70$. At the selected operating points the output matches the baseline whenever no conflict is present, so the contrastive pass is conditional rather than mandatory.

\Figure[t!p](topskip=0pt, botskip=0pt, midskip=0pt)[width=\textwidth]{fig6.png}{The same sweep for the three naturally balanced models, where EGCD recovers the greedy baseline at sufficiently high thresholds for the two LLaVA models while Universal CD+APC falls up to 10.62 points below it.\label{fig:sound}}

\subsection{First-Token Entropy and Gating Threshold Dynamics}\label{subsec:dynamics}

The sweep shows what the routing decision is actually doing. For Qwen2-VL-2B, balanced accuracy rises to its peak at the selected threshold (held-out numbers in Table~\ref{tab:main}) as lower-entropy interventions are excluded from System 2 (Figure~\ref{fig:sweep}). Past 0.70 the curve slides back toward the greedy baseline. A gate set too high lets real conflicts slip through. The response is hill-shaped, and where the peak sits depends on each model's entropy distribution. That is why no single $\tau$ works for all seven models.

The naturally balanced models show the flat side of the same pattern. Idefics2-8B stays inside a 0.7-point band across the whole sweep (Figure~\ref{fig:sound}), which means its first-token entropy stays below the threshold during normal operation, so the costlier contrastive pass runs only when entropy crosses it. Figure~\ref{fig:all} gives the full sweep for all seven models, and the tuning-half numbers are in Appendix~\ref{app:sweep}. The two views together describe a gate that gets the cost of the intervention right in both directions. It pays for contrast where conflict is present and withholds it where there is none, which keeps inference cost in line.

\Figure[t!p](topskip=0pt, botskip=0pt, midskip=0pt)[width=\textwidth]{fig7.png}{Gain over greedy decoding across the entropy-threshold sweep for all seven models: every curve falls back toward zero as $\tau$ grows, and on the naturally balanced models a loose threshold costs up to 10.62 points.\label{fig:all}}

On the held-out split, the first-token entropy predicts per-item correctness with a binary error-prediction AUC between 0.62 (SmolVLM2-2.2B) and 0.79 (Qwen2.5-VL-7B), with the remaining five models between 0.70 and 0.74 (Qwen2-VL-2B 0.71, Phi-3.5-Vision-4.2B 0.70, Idefics2-8B 0.74, LLaVA-1.5-7B 0.73, LLaVA-v1.6-Mistral-7B 0.73). The entropy probe thus carries signal about which queries the model answers incorrectly; whether high-entropy queries specifically benefit from CD, stratified by entropy and label, is left for future work.

The gate routes 40.9\% of held-out queries on Qwen2-VL-2B, 48.3\% on Phi-3.5-Vision-4.2B, 28.0\% on SmolVLM2-2.2B, 14.6\% on Qwen2.5-VL-7B, 6.7\% on Idefics2-8B, none on LLaVA-1.5-7B, and 0.2\% on LLaVA-v1.6-Mistral-7B. Shuffled routing, a control that sends the same fraction of queries to System 2 as the gate but selects them uniformly at random instead of by entropy, recovers only part of the gain (64.50 against 67.33 on SmolVLM2-2.2B, 73.41 against 80.45 on Qwen2-VL-2B, 72.79 against 77.11 on Phi-3.5-Vision-4.2B), which suggests that most of the improvement comes from which queries are routed rather than from how often the intervention fires.

\section{Conclusion and Future Work}\label{sec:conclusion}

\subsection{Summary of Contributions}\label{subsec:summary}

This paper treats relational hallucination as a decoding failure with a hypothesised mechanism behind it, Visual Inertia: cross-attention settles on its first targets, grounding decays with every token generated, and the language priors end up writing the answer \cite{gong2026attention}. EGCD addresses it with a query-level entropy gate that fires a counterfactual contrastive pass only where entropy crosses a calibrated threshold, bounded by an Adaptive Plausibility Constraint and decoded greedily. The link from the first-token entropy probe to the underlying attention mechanism is treated as a hypothesis, supported by the error-prediction AUC reported in Section~\ref{subsec:dynamics}.

Across the seven models on AMBER's relation subset, the gate separated the suite without being told which was which. The biased architectures gained, Qwen2-VL-2B by 9.74 points and Phi-3.5-Vision-4.2B by 8.06, with a smaller specificity cost than universal CD+APC. The naturally balanced models kept their numbers: universal CD had pushed LLaVA-1.5-7B down to 65.09\%, while EGCD handed back its 75.11\% baseline essentially unchanged. A shared suite-wide threshold keeps most of the gains on the biased models but gives up the idle-gate safety on LLaVA-1.5-7B: forcing the same $\tau = 0.69$ on every model routes 23\% of its held-out queries to System 2 and drops balanced accuracy to 72.97 against the 75.11 baseline. The result is a conditional, per-model calibrated use of contrastive decoding that spends its correction only on the queries that need it.

\subsection{Limitations and Future Work}\label{subsec:future}

Several extensions look worth the effort. The first is calibration-free gating: entropy percentiles estimated on the fly, or one shared threshold fitted across a model family, would remove the per-model sweep that $\tau^{*}$ currently requires. Open-ended generation is the second and harder target, since the first-token distribution no longer concentrates the whole decision in one place and some form of step-level probe would be needed. The third is routing inside reasoning models: Octopus learns per-step strategy selection \cite{suo2025octopus}, dual-process search keeps spreading through chain-of-thought systems \cite{cheng2025think,li2025system}, and a training-free entropy signal is a natural candidate for deciding when those pipelines should slow down and think twice. Two further extensions follow directly from the methodology: (i) measure attention trajectories on the evaluated models to test the Visual Inertia hypothesis directly, rather than relying on the first-token entropy proxy; (ii) run a stratified analysis of CD fixes versus harms by entropy and label to verify that high-entropy queries specifically benefit from the contrastive pass.

\section*{Acknowledgment}

The authors thank the anonymous reviewers for constructive feedback that improved the presentation of this work.

\section*{Declarations}

\begin{itemize}
\item \textbf{Funding:} Not applicable.
\item \textbf{Conflict of interest/Competing interests:} The authors declare no competing interests.
\item \textbf{Ethics approval and consent to participate:} Not applicable.
\item \textbf{Consent for publication:} All authors consent to the publication of this manuscript.
\item \textbf{Data availability:} The AMBER benchmark used in this study is publicly available at \url{https://github.com/junyangwang0410/AMBER}. The canonical split (822 tuning / 842 held-out query ids, seed = 42, stratified by image id) is released alongside the code.
\item \textbf{Materials availability:} Not applicable.
\item \textbf{Code availability:} The EGCD inference pipeline, the evaluation scripts, the per-item correctness outputs used for the exact paired McNemar tests, and the split IDs will be made available upon publication.
\item \textbf{Author contribution:} All authors contributed equally to the design of the methodology, the implementation of the experiments, and the writing of the manuscript. All authors read and approved the final manuscript.
\end{itemize}

\appendices
\section{Complete Entropy-Threshold Sweep Results}\label{app:sweep}

Table~\ref{tab:sweep} reports the balanced accuracy of the entropy-threshold sweep on the tuning half (822 queries). It is the selection record: each model's peak here is the operating point $\tau^{*}$ that Section~\ref{sec:results} evaluates on the held-out half. The selected thresholds are 0.70 (Qwen2-VL-2B), 0.57 (Phi-3.5-Vision-4.2B), 0.68 (SmolVLM2-2.2B), 0.64 (Qwen2.5-VL-7B), 0.66 (Idefics2-8B), 0.70 (LLaVA-1.5-7B), and 0.70 (LLaVA-v1.6-Mistral-7B); recall and specificity at each operating point appear in Table~\ref{tab:main}. Ties in the tuning-half argmax are broken by selecting the lowest $\tau$. Figure~\ref{fig:heatmap} renders the held-out results of Table~\ref{tab:main} as three heatmaps and Figure~\ref{fig:gain} summarises the gain each selected threshold buys.

\begin{table*}[!tp]
\caption{Balanced accuracy (\%) across the entropy-threshold sweep on the tuning half for all seven models; each model's peak on this grid is its selected operating point $\tau^{*}$, evaluated on the held-out half in Table~\ref{tab:main}.}\label{tab:sweep}
\begin{tabular*}{\textwidth}{@{\extracolsep\fill}lccccccc@{}}
\toprule
$\tau$ & \textbf{Qwen2-VL-2B} & \textbf{Phi-3.5-Vision-4.2B} & \textbf{SmolVLM2-2.2B} & \textbf{Qwen2.5-VL-7B} & \textbf{Idefics2-8B} & \textbf{LLaVA-1.5-7B} & \textbf{LLaVA-v1.6-Mistral-7B}\\
\midrule
0.55 & 76.34 & 76.74 & 67.59 & 81.82 & 82.80 & 68.00 & 77.99 \\
0.56 & 76.34 & 76.79 & 67.59 & 81.82 & 82.80 & 68.00 & 78.09 \\
0.57 & 76.34 & \textbf{77.14} & 67.59 & 81.82 & 82.80 & 68.00 & 78.09 \\
0.58 & 76.34 & 77.14 & 67.59 & 81.92 & 82.80 & 68.00 & 78.39 \\
0.59 & 76.34 & 76.78 & 67.59 & 81.82 & 82.80 & 68.00 & 78.80 \\
0.60 & 76.49 & 76.98 & 67.59 & 81.82 & 82.80 & 68.10 & 79.05 \\
0.61 & 76.49 & 77.08 & 67.59 & 81.87 & 82.80 & 68.41 & 79.56 \\
0.62 & 76.65 & 77.03 & 67.59 & 81.87 & 82.80 & 68.71 & 79.86 \\
0.63 & 76.80 & 76.68 & 67.59 & 82.23 & 82.80 & 68.81 & 80.62 \\
0.64 & 76.80 & 76.42 & 67.74 & \textbf{82.38} & 82.80 & 69.22 & 81.43 \\
0.65 & 77.71 & 75.86 & 67.74 & 82.38 & 82.80 & 70.13 & 82.04 \\
0.66 & 77.86 & 75.81 & 67.74 & 82.38 & \textbf{83.11} & 70.49 & 82.35 \\
0.67 & 77.91 & 75.85 & 67.74 & 82.17 & 83.11 & 71.05 & 82.25 \\
0.68 & 77.50 & 74.68 & \textbf{68.19} & 81.77 & 83.00 & 71.97 & 83.46 \\
0.69 & 78.01 & 73.56 & 67.32 & 81.57 & 83.00 & 75.08 & 84.48 \\
0.70 & \textbf{78.31} & 72.24 & 66.34 & 81.82 & 83.00 & \textbf{78.05} & \textbf{86.40} \\
0.71 & 77.29 & 69.95 & 65.06 & 81.61 & 83.00 & 78.05 & 86.40 \\
0.72 & 75.86 & 69.69 & 64.76 & 80.85 & 83.00 & 78.05 & 86.40 \\
0.73 & 74.88 & 69.69 & 64.76 & 80.39 & 83.00 & 78.05 & 86.40 \\
0.74 & 72.73 & 69.59 & 64.76 & 80.65 & 83.00 & 78.05 & 86.40 \\
0.75 & 72.43 & 69.59 & 64.76 & 80.55 & 83.00 & 78.05 & 86.40 \\
\bottomrule
\end{tabular*}
\end{table*}

\Figure[t!p](topskip=0pt, botskip=0pt, midskip=0pt)[width=\textwidth]{fig8.png}{The held-out results of Table~\ref{tab:main} rendered as three heatmaps (Recall, Specificity, Balanced Accuracy; seven models by five decoding regimes), with the EGCD column outlined.\label{fig:heatmap}}

\Figure[t!p](topskip=0pt, botskip=0pt, midskip=0pt)[width=\textwidth]{fig9.png}{Gain over greedy decoding at each model's selected threshold $\tau^{*}$ on the tuning half: the gate buys the most exactly where hallucination bias is largest.\label{fig:gain}}

\clearpage

\bibliographystyle{IEEEtran}
\bibliography{references}

\EOD

\end{document}
